\documentclass{article}
\usepackage[utf8]{inputenc}
\usepackage{amsmath,amssymb}
\usepackage[most]{tcolorbox}
\usepackage{geometry}
\geometry{margin=2.5cm}

\newcommand{\step}[1]{\par\noindent\hangindent=3.5em\hangafter=1%
  \makebox[3.5em][l]{\textbf{#1.}}}
\newcommand{\steptag}[1]{\hfill{\small\textsf{[#1]}}}

\newtcolorbox{proofbox}[1]{
  colback=gray!5, colframe=gray!60,
  fonttitle=\bfseries, title={#1},
  breakable, enhanced,
  left=4pt, right=4pt, top=4pt, bottom=4pt,
  before skip=10pt, after skip=10pt
}

\begin{document}

\begin{proofbox}{Top concentration: for an exact signed idempotent P with ...}
\step{1} Top concentration: for an exact signed idempotent P with 0 < delta(P) <= 1/4, nonempty visible set W(P), hidden top vertex v of height H, and halo width a > 0 with H > a*tau (tau = sqrt(delta)), the positive mass v places outside the genuine set satisfies sum\_\{j notin G\_a\} P\_vj\textasciicircum{}+ <= nu\_v*(2+4*delta)/(H - a*tau), where G\_a = \{j : dist\_1(p\_j, conv\{p\_w : w in W\}) > a*tau\} and nu\_v is the row-v negative mass. \steptag{validated} \\[6pt]
\step{1.1} Set C=conv\{p\_w:w in W\}, d\_j=dist\_1(p\_j,C), D=2+4*delta, and choose an l1/l\_infty support functional phi with ||u||\_inf<=1, phi<=0 on C, phi(p\_v)=H; then, for t\_j:=H-phi(p\_j), every row satisfies 0<=t\_j<=D and every j notin G\_a satisfies t\_j>=H-a*tau. \steptag{validated} \\[6pt]
\step{1.1.1} Because W is nonempty, C=conv\{p\_w:w in W\} is a nonempty compact convex subset of finite-dimensional l1 space. Since v is a hidden top vertex of height H, H=dist\_1(p\_v,C). The l1/l\_infty distance duality formula gives u with ||u||\_inf<=1 such that u.p\_v - sup\_\{c in C\} u.c = H; with b=-sup\_C u.c and phi(x)=u.x+b, we have phi<=0 on C and phi(p\_v)=H. \steptag{validated} \\[4pt]
\step{1.1.2} For every row p\_j, phi(p\_j)<=d\_j:=dist\_1(p\_j,C): indeed phi(c)<=0 for c in C and ||u||\_inf<=1 give phi(p\_j)<=u.(p\_j-c)<=||p\_j-c||\_1, then inf over c in C. Since H is the maximum row height, d\_j<=H for all j; if j notin G\_a, then d\_j<=a*tau. \steptag{validated} \\[4pt]
\step{1.1.3} The functional phi(x)=u.x+b has ||u||\_inf<=1, hence is 1-Lipschitz for the l1 metric. By the row-geometry clause of def-signed-idempotent, ||p\_v-p\_j||\_1<=D:=2+4*delta for every row p\_j, so t\_j:=H-phi(p\_j)=phi(p\_v)-phi(p\_j)<=D. Also phi<=0 on C and ||u||\_inf<=1 imply phi(p\_j)<=d\_j:=dist\_1(p\_j,C); since v is a hidden top vertex of height H, d\_j<=H for all rows, and by the definition of G\_a, j notin G\_a gives d\_j<=a*tau. Therefore t\_j>=H-d\_j>=0 for all j and t\_j>=H-a*tau for j notin G\_a. \steptag{validated} \\[4pt]
\step{1.2} Dependency-scoped affine reproduction: for the support functional phi supplied by node 1.1 (so phi is affine and phi(p\_v)=H), setting t\_j:=H-phi(p\_j), exact idempotence gives sum\_j P\_vj*t\_j=0. \steptag{validated} \\[6pt]
\step{1.2.1} From def-signed-idempotent, P*1=1 and P\textasciicircum{}2=P. The row-v equation of P\textasciicircum{}2=P is p\_v=sum\_j P\_vj*p\_j as a vector, and P*1=1 gives sum\_j P\_vj=1. \steptag{validated} \\[4pt]
\step{1.2.2} For any affine phi, affine-linearity over coefficients with total sum 1 gives phi(p\_v)=sum\_j P\_vj*phi(p\_j). Using phi(p\_v)=H and sum\_j P\_vj=1, we get sum\_j P\_vj*(H-phi(p\_j))=H*sum\_j P\_vj-sum\_j P\_vj*phi(p\_j)=H-H=0. \steptag{validated} \\[4pt]
\step{1.2.3} Scoped bridge for challenge ch-68153d29c48e7b18: node 1.1 supplies the chosen affine support functional phi with phi(p\_v)=H, and nodes 1.2.1 and 1.2.2 supply the exact-idempotence row equation, total-mass equation, and affine reproduction calculation. Therefore, for this phi and t\_j:=H-phi(p\_j), sum\_j P\_vj*t\_j = H*sum\_j P\_vj - sum\_j P\_vj*phi(p\_j) = H - phi(p\_v) = H - H = 0. This is exactly the amended dependency-scoped statement of node 1.2. \steptag{validated} \\[4pt]
\step{1.3} The nonnegative bounded deficits t\_j obey the sign-split estimate sum\_j P\_vj\textasciicircum{}+*t\_j <= nu\_v*(2+4*delta). \steptag{validated} \\[6pt]
\step{1.3.1} Write P\_vj\textasciicircum{}+=max(P\_vj,0) and P\_vj\textasciicircum{}-=max(-P\_vj,0), so P\_vj=P\_vj\textasciicircum{}+-P\_vj\textasciicircum{}- and nu\_v=sum\_j P\_vj\textasciicircum{}- by the row-v negative-mass notation from def-negative-mass and the contract. Establish locally that t\_j satisfies 0<=t\_j<=D:=2+4*delta for every j and that sum\_j P\_vj*t\_j=0. Then sum\_j P\_vj\textasciicircum{}+*t\_j=sum\_j P\_vj\textasciicircum{}-*t\_j, and this common value is at most D*sum\_j P\_vj\textasciicircum{}-=nu\_v*(2+4*delta). \steptag{validated} \\[6pt]
\step{1.3.1.1} Local bounded-deficit premise: let C=conv\{p\_w:w in W\}, let phi(x)=u.x+b be the support functional used for the height H with ||u||\_inf<=1, phi<=0 on C, and phi(p\_v)=H, and set t\_j:=H-phi(p\_j), D:=2+4*delta. For every row p\_j, phi(p\_j)<=d\_j:=dist\_1(p\_j,C)<=H, so t\_j>=0. Also phi is 1-Lipschitz for the l1 metric and def-signed-idempotent gives ||p\_v-p\_j||\_1<=D for every row, hence t\_j=phi(p\_v)-phi(p\_j)<=D. Therefore 0<=t\_j<=D for every j. \steptag{validated} \\[4pt]
\step{1.3.1.2} Local affine-reproduction premise: def-signed-idempotent gives P*1=1 and P\textasciicircum{}2=P. Taking row v in P\textasciicircum{}2=P gives p\_v=sum\_j P\_vj*p\_j as a vector, and P*1=1 gives sum\_j P\_vj=1. For the affine phi(x)=u.x+b, sum\_j P\_vj*phi(p\_j)=u.(sum\_j P\_vj*p\_j)+b*sum\_j P\_vj=u.p\_v+b=phi(p\_v)=H. Consequently sum\_j P\_vj*t\_j=sum\_j P\_vj*(H-phi(p\_j))=H*sum\_j P\_vj-sum\_j P\_vj*phi(p\_j)=H-H=0. \steptag{validated} \\[4pt]
\step{1.3.1.3} Using 1.3.1.1 and 1.3.1.2, write P\_vj=P\_vj\textasciicircum{}+-P\_vj\textasciicircum{}- with P\_vj\textasciicircum{}+,P\_vj\textasciicircum{}- >=0 and nu\_v=sum\_j P\_vj\textasciicircum{}- by the row-v negative-mass notation. Then 0=sum\_j P\_vj*t\_j=sum\_j P\_vj\textasciicircum{}+*t\_j-sum\_j P\_vj\textasciicircum{}-*t\_j, so sum\_j P\_vj\textasciicircum{}+*t\_j=sum\_j P\_vj\textasciicircum{}-*t\_j. Since 0<=t\_j<=D and P\_vj\textasciicircum{}- >=0 for every j, sum\_j P\_vj\textasciicircum{}-*t\_j<=D*sum\_j P\_vj\textasciicircum{}-=nu\_v*(2+4*delta). Thus sum\_j P\_vj\textasciicircum{}+*t\_j<=nu\_v*(2+4*delta). \steptag{validated} \\[4pt]
\step{1.4} Since j notin G\_a implies t\_j>=H-a*tau>0, the sign-split estimate gives (H-a*tau)*sum\_\{j notin G\_a\} P\_vj\textasciicircum{}+ <= nu\_v*(2+4*delta), and division by H-a*tau proves the claimed bound. \steptag{validated} \\[6pt]
\step{1.4.1} In the local setting of node 1.4, choose the l1/l\_infty support functional phi for the distance from p\_v to C=conv\{p\_w:w in W\}, put t\_j=H-phi(p\_j), and put D=2+4*delta. The support-functional bounds give t\_j>=0 for all j and t\_j>=H-a*tau for j notin G\_a; exact idempotence and the row-v negative-mass definition give the sign-split estimate sum\_j P\_vj\textasciicircum{}+*t\_j<=nu\_v*D. Hence (H-a*tau)*sum\_\{j notin G\_a\} P\_vj\textasciicircum{}+ <= sum\_\{j notin G\_a\} P\_vj\textasciicircum{}+*t\_j <= sum\_j P\_vj\textasciicircum{}+*t\_j <= nu\_v*(2+4*delta). Since H>a*tau, division by H-a*tau gives sum\_\{j notin G\_a\} P\_vj\textasciicircum{}+ <= nu\_v*(2+4*delta)/(H-a*tau). \steptag{validated} \\[6pt]
\step{1.4.1.1} Dependency bridge for the final division step: the only nonlocal inputs used here are exactly node 1.1, which gives t\_j>=0 for all j and t\_j>=H-a*tau for j notin G\_a, and node 1.3, which gives sum\_j P\_vj\textasciicircum{}+*t\_j<=nu\_v*(2+4*delta). With these two dependencies explicitly declared, the rest of node 1.4.1 is pure order algebra: because P\_vj\textasciicircum{}+=max(P\_vj,0)>=0, multiplying the lower bound on t\_j over j notin G\_a gives (H-a*tau)*sum\_\{j notin G\_a\}P\_vj\textasciicircum{}+<=sum\_\{j notin G\_a\}P\_vj\textasciicircum{}+*t\_j; since t\_j>=0 and P\_vj\textasciicircum{}+>=0, the restricted sum is at most the full sum; node 1.3 bounds the full sum by nu\_v*(2+4*delta); and H>a*tau makes H-a*tau>0, so division by this positive scalar gives the desired bound. \steptag{archived} \\[4pt]
\step{1.4.1.2} Scoped dependency bridge for node 1.4.1: assume exactly the explicitly declared dependencies 1.1 and 1.3. Dependency 1.1 supplies t\_j>=0 for every j and t\_j>=H-a*tau for every j notin G\_a; dependency 1.3 supplies sum\_j P\_vj\textasciicircum{}+*t\_j<=nu\_v*(2+4*delta). Since P\_vj\textasciicircum{}+=max(P\_vj,0)>=0, summing the lower bound over j notin G\_a gives (H-a*tau)*sum\_\{j notin G\_a\} P\_vj\textasciicircum{}+ <= sum\_\{j notin G\_a\} P\_vj\textasciicircum{}+*t\_j. Because t\_j>=0 and P\_vj\textasciicircum{}+>=0, the restricted sum is at most sum\_j P\_vj\textasciicircum{}+*t\_j. Applying 1.3 gives (H-a*tau)*sum\_\{j notin G\_a\} P\_vj\textasciicircum{}+ <= nu\_v*(2+4*delta). Finally H>a*tau makes H-a*tau>0, so division by H-a*tau yields sum\_\{j notin G\_a\} P\_vj\textasciicircum{}+ <= nu\_v*(2+4*delta)/(H-a*tau). \steptag{validated} \\[4pt]
\step{1.4.1.3} Self-contained dependency repair for ch-0e785a8c7af38a9a: no sibling node is used. Let C=conv\{p\_w:w in W\} and d\_j=dist\_1(p\_j,C). Since W is nonempty and v is a hidden top vertex of height H=dist\_1(p\_v,C), finite-dimensional l1/l\_infty distance duality supplies an affine phi(x)=u.x+b with ||u||\_inf<=1, phi<=0 on C, and phi(p\_v)=H. Set t\_j=H-phi(p\_j) and D=2+4*delta. For every row p\_j, phi(p\_j)<=d\_j: for c in C, phi(p\_j)<=phi(p\_j)-phi(c)=u.(p\_j-c)<=||p\_j-c||\_1, then take the infimum over c. Since H is the maximum row height, d\_j<=H, so t\_j>=0. If j notin G\_a, then d\_j<=a*tau by the definition G\_a=\{j:dist\_1(p\_j,C)>a*tau\}, hence t\_j>=H-d\_j>=H-a*tau. Also phi is 1-Lipschitz for l1 and def-signed-idempotent gives ||p\_v-p\_j||\_1<=D for every row, so t\_j=phi(p\_v)-phi(p\_j)<=D. Exact signed idempotence gives P*1=1 and P\textasciicircum{}2=P, hence sum\_j P\_vj=1 and p\_v=sum\_j P\_vj p\_j. Affineness gives sum\_j P\_vj phi(p\_j)=phi(p\_v)=H, so sum\_j P\_vj t\_j=0. Writing P\_vj=P\_vj\textasciicircum{}+-P\_vj\textasciicircum{}- with P\_vj\textasciicircum{}+,P\_vj\textasciicircum{}->=0 and nu\_v=sum\_j P\_vj\textasciicircum{}- by the row-v negative-mass notation, the equality sum\_j P\_vj t\_j=0 gives sum\_j P\_vj\textasciicircum{}+ t\_j=sum\_j P\_vj\textasciicircum{}- t\_j<=D*sum\_j P\_vj\textasciicircum{}-=nu\_v*(2+4*delta). Combining this with P\_vj\textasciicircum{}+>=0, t\_j>=0, and t\_j>=H-a*tau on j notin G\_a yields (H-a*tau)*sum\_\{j notin G\_a\} P\_vj\textasciicircum{}+<=sum\_\{j notin G\_a\} P\_vj\textasciicircum{}+ t\_j<=sum\_j P\_vj\textasciicircum{}+ t\_j<=nu\_v*(2+4*delta). Since H>a*tau, divide by the positive number H-a*tau to obtain the asserted bound. \steptag{validated} \\[4pt]
\end{proofbox}

\end{document}
